\documentclass[12pt, openany]{book}

% ──────────────────────────────────────────────
% Packages
% ──────────────────────────────────────────────
\usepackage[margin=1in]{geometry}
\usepackage{amsmath, amssymb, amsthm}
\usepackage{mathtools}
\usepackage{bm}
\usepackage{tikz}
\usepackage{pgfplots}
\pgfplotsset{compat=1.18}
\usetikzlibrary{arrows.meta, decorations.markings, calc, patterns, positioning}
\usepackage{tcolorbox}
\tcbuselibrary{theorems, breakable, skins}
\usepackage{hyperref}
\usepackage{enumitem}
\usepackage{fancyhdr}
\usepackage{titlesec}
\usepackage{epigraph}
\usepackage{xcolor}
\usepackage{booktabs}

% ──────────────────────────────────────────────
% Colors
% ──────────────────────────────────────────────
\definecolor{chapblue}{RGB}{0, 70, 130}
\definecolor{accentred}{RGB}{180, 40, 40}
\definecolor{softgray}{RGB}{245, 245, 245}
\definecolor{trajgreen}{RGB}{30, 130, 60}
\definecolor{darkgold}{RGB}{160, 120, 20}
\definecolor{purplehaze}{RGB}{100, 40, 140}

% ──────────────────────────────────────────────
% Theorem environments
% ──────────────────────────────────────────────
\newtcbtheorem[number within=chapter]{principle}{Principle}{
  colback=chapblue!5, colframe=chapblue!80,
  fonttitle=\bfseries, separator sign={.~},
  breakable
}{princ}

\newtcbtheorem[number within=chapter]{keyidea}{Key Idea}{
  colback=accentred!5, colframe=accentred!70,
  fonttitle=\bfseries, separator sign={.~},
  breakable
}{idea}

\newtcbtheorem[number within=chapter]{example}{Example}{
  colback=softgray, colframe=black!40,
  fonttitle=\bfseries, separator sign={.~},
  breakable
}{ex}

\newtcbtheorem[number within=chapter]{definition}{Definition}{
  colback=darkgold!5, colframe=darkgold!80,
  fonttitle=\bfseries, separator sign={.~},
  breakable
}{def}

\newtcbtheorem[number within=chapter]{remark}{Remark}{
  colback=purplehaze!5, colframe=purplehaze!60,
  fonttitle=\bfseries, separator sign={.~},
  breakable
}{rem}

\newtcbtheorem[number within=chapter]{warning}{Warning}{
  colback=accentred!5, colframe=accentred!80,
  fonttitle=\bfseries, separator sign={.~},
  breakable
}{warn}

\theoremstyle{plain}
\newtheorem{proposition}{Proposition}[chapter]
\newtheorem{theorem}{Theorem}[chapter]
\newtheorem{lemma}[theorem]{Lemma}
\newtheorem{corollary}[theorem]{Corollary}

% ──────────────────────────────────────────────
% Chapter formatting
% ──────────────────────────────────────────────
\titleformat{\chapter}[display]
  {\normalfont\huge\bfseries\color{chapblue}}
  {\chaptertitlename\ \thechapter}{20pt}{\Huge}
\titlespacing*{\chapter}{0pt}{-20pt}{30pt}

% ──────────────────────────────────────────────
% Header/Footer
% ──────────────────────────────────────────────
\pagestyle{fancy}
\fancyhf{}
\fancyhead[LE,RO]{\thepage}
\fancyhead[RE]{\textit{Control Is Motion}}
\fancyhead[LO]{\textit{\nouppercase{\leftmark}}}
\renewcommand{\headrulewidth}{0.4pt}
\setlength{\headheight}{14.5pt}

% ──────────────────────────────────────────────
% Custom commands
% ──────────────────────────────────────────────
\newcommand{\state}{\bm{x}}
\newcommand{\control}{\bm{u}}
\newcommand{\traj}{\bm{\gamma}}
\newcommand{\manifold}{\mathcal{M}}
\newcommand{\configspace}{\mathcal{Q}}
\newcommand{\statespace}{\mathcal{X}}
\newcommand{\controlspace}{\mathcal{U}}
\newcommand{\Reals}{\mathbb{R}}
\newcommand{\dd}{\mathrm{d}}
\newcommand{\dt}{\,\dd t}
\newcommand{\norm}[1]{\left\lVert #1 \right\rVert}
\newcommand{\inner}[2]{\left\langle #1,\, #2 \right\rangle}

% ──────────────────────────────────────────────
\begin{document}

\setcounter{chapter}{6}

\chapter{Funnel Synthesis}

\epigraph{%
  You cannot control the exact trajectory a stone will take down a mountain.\\
  But you can build a trough to ensure it reaches the bottom.
}{}

\vspace{0.5cm}

% ══════════════════════════════════════════════
\section{Beyond LQR: The Need for Guaranteed Invariance}
% ══════════════════════════════════════════════

In Chapter 4, we showed that applying time-varying Linear Quadratic Regulation (LQR) to the transverse dynamics creates a contracting envelope of stability. The Riccati equation naturally shapes this envelope into a "funnel" that narrows as the terminal constraint approaches. 

However, LQR guarantees are inherently \emph{local}. They rely on the linear approximation of the transverse dynamics $\dot{\bm{\xi}} = \bm{A}_\perp(t)\bm{\xi} + \bm{B}_\perp(t)\delta\control$. For actual physical systems traversing highly nonlinear state spaces—like a golf club accelerating to $50$ m/s while subject to massive centrifugal and Coriolis forces—this linearization quickly breaks down. If a disturbance pushes the state too far from the reference trajectory $\traj^*(t)$, the nonlinearities will dominate, the LQR restoring forces will compute the wrong corrective action, and the clubhead will completely miss the ball.

We need a mathematically rigorous way to define the \emph{exact boundary} of our funnel. We must prove that for the fully nonlinear equations of motion, any state starting within the mouth of the funnel at time $t=0$ will remain inside the funnel for all $t \in [0, t_f]$.

\begin{definition}{Dynamic Funnel}{dynamic_funnel}
  Given a reference trajectory $\traj^*(t)$, a feedback controller $\control = \bm{k}(\state, t)$, and a scalar boundary function $V(\state, t)$, the set:
  \begin{equation}
    \mathcal{F}(t) = \{ \state \in \statespace \mid V(\state, t) \leq 1 \}
  \end{equation}
  constitutes a valid \textbf{dynamic funnel} if the set is positively invariant under the closed-loop dynamics. That is, if $\state(0) \in \mathcal{F}(0)$, then $\state(t) \in \mathcal{F}(t)$ for all $t > 0$.
\end{definition}

% ══════════════════════════════════════════════
\section{Time-Varying Lyapunov Functions for Trajectories}
% ══════════════════════════════════════════════

To verify invariance, we use the machinery of Lyapunov theory, adapted for trajectories rather than equilibria. We require our boundary function $V(\state, t)$ to act as a time-varying Lyapunov function.

\begin{theorem}[Funnel Invariance]\label{thm:funnel_invariance}
  The set $\mathcal{F}(t) = \{ \state \mid V(\state, t) \leq 1 \}$ is invariant if, on the boundary where $V(\state, t) = 1$, the function is strictly decreasing along the system trajectories:
  \begin{equation}
    \dot{V}(\state, t) = \frac{\partial V}{\partial t} + \nabla_{\state} V(\state, t) \cdot \bm{f}(\state, \bm{k}(\state, t)) < 0 \quad \text{for all } \state \text{ s.t. } V(\state, t) = 1.
  \end{equation}
\end{theorem}

If the derivative of the boundary function points inward everywhere on the boundary, trajectories may circulate inside the funnel, but they can never escape it. For a moving target problem, we want the spatial volume of $\mathcal{F}(t)$ to shrink as $t \to t_f$, terminating in a tight bounded region at impact.

\begin{keyidea}{Funnels Define the Usable State Space}{usable_space}
  A computed funnel is not just an analysis tool; it is the ultimate judge of repeatability. It explicitly defines the exact set of initial conditions (e.g., the kinematic state at the top of the backswing) that are mathematically guaranteed to successfully strike the ball despite the presence of extreme nonlinearities.
\end{keyidea}

% ══════════════════════════════════════════════
\section{Sums-of-Squares (SOS) Programming}
% ══════════════════════════════════════════════

Proving that $\dot{V} < 0$ on the boundary $V=1$ for a general nonlinear function $\bm{f}$ is famously difficult. We have an infinite number of states on the boundary to check! We circumvent this by restricting our dynamics $\bm{f}$ algorithms to polynomials (via Taylor expansion or substitution) and leveraging \textbf{Sums-of-Squares (SOS) Programming}.

A polynomial $p(x)$ is a Sum of Squares if it can be written as $p(x) = \sum_{i} q_i(x)^2$ for some polynomials $q_i(x)$. If a polynomial is SOS, it is globally non-negative. Testing if a polynomial is SOS turns out to be equivalent to solving a Semidefinite Program (SDP)—a convex optimization problem that can be solved very efficiently.

We can reframe Theorem \ref{thm:funnel_invariance} using the S-procedure. To guarantee that $\dot{V}(\state, t) < 0$ whenever $V(\state, t) = 1$, we require:
\begin{equation}
  -\dot{V}(\state, t) - L(\state, t)(1 - V(\state, t)) \quad \text{is SOS}
\end{equation}
where $L(\state, t)$ is a strictly positive polynomial multiplier. If $V(\state, t) = 1$, the second term vanishes, and the SOS condition forces $-\dot{V} > 0$, achieving our invariance proof!

% ══════════════════════════════════════════════
\section{Computing the Funnel for the Golf Swing}
% ══════════════════════════════════════════════

Let us return to the double pendulum of the golf swing. Our reference trajectory $\traj^*(t)$ from Chapter 6 gives us the nominal kinematics. Our transverse LQR from Chapter 4 gives us a feedback matrix $\bm{K}(t)$ designed to stabilize deviations. We wish to compute the largest possible funnel around this sequence that the closed-loop system is mathematically guaranteed to stay inside.

We define our candidate Lyapunov function as a time-varying quadratic form:
\begin{equation}
  V(\state, t) = (\state - \traj^*(t))^\top \bm{P}(t) (\state - \traj^*(t))
\end{equation}
where $\bm{P}(t)$ is a positive definite matrix that changes continuously along the trajectory. The funnel boundary is the time-varying ellipsoid where $V=1$. By altering $\bm{P}(t)$, we alter the shape and orientation of the funnel.

\begin{example}{Volume Maximization via SDP}{vol_max}
  We transcribe the continuous time $t$ into discrete nodes $t_k$, just as we did in direct collocation. At each node, we set up an SDP that seeks to \emph{maximize} the volume of the ellipsoid defined by $\bm{P}(t_k)$, subject to the SOS invariance constraints connecting it to adjacent nodes.

  The optimizer actively computes the exact structural limits of the golf swing's passive dynamics. It discovers that early in the downswing, the funnel can be incredibly "wide" (forgiving), but as centrifugal forces build, the SOS constraints force $\bm{P}(t_k)$ to become highly elongated along the unactuated dimensions in order to maintain the invariance guarantee. The resulting geometry is a true numerical portrait of athletic consistency.
\end{example}

In Chapter 8, we will explore how to make these funnels robust to temporal variation by migrating out of the time domain entirely and into the phase domain.

\end{document}
